\documentclass[11pt]{article}

\usepackage[margin=0.85in]{geometry}
\usepackage{amsmath,amssymb}
\usepackage{enumitem}
\usepackage{xcolor}
\usepackage{tcolorbox}
\usepackage{fancyhdr}
\usepackage{lastpage}
\usepackage{hyperref}
\usepackage{microtype}

\definecolor{uwpurple}{HTML}{4B2E83}
\definecolor{uwgold}{HTML}{B7A57A}
\definecolor{softpurple}{HTML}{F4F1FA}

\hypersetup{
    colorlinks=true,
    linkcolor=uwpurple,
    urlcolor=uwpurple
}

\pagestyle{fancy}
\fancyhf{}
\lhead{\textbf{MATH 394: Probability I}}
\rhead{\textbf{Homework 2}}
\cfoot{\thepage\ of \pageref{LastPage}}

\setlength{\parindent}{0pt}
\setlength{\parskip}{0.45em}

\newcommand{\course}{MATH 394: Probability I}
\newcommand{\instructor}{Arman Jahangiri}
\newcommand{\department}{Department of Mathematics}
\newcommand{\university}{University of Washington}

\newtcolorbox{problemheader}{
    colback=softpurple,
    colframe=uwpurple,
    boxrule=0.8pt,
    arc=2mm,
    left=2mm,
    right=2mm,
    top=1mm,
    bottom=1mm
}

\newcommand{\source}[1]{%
    \vspace{0.15cm}
    {\small\textit{Source: Blitzstein and Hwang, \textit{Introduction to Probability}, Chapter 1, Problem #1.}}
}

\newcommand{\answerbox}{%
\begin{flushright}
\begin{tcolorbox}[
    width=0.36\textwidth,
    height=2cm,
    colback=white,
    colframe=uwpurple,
    boxrule=0.7pt,
    arc=1mm
]
{\small\textbf{Final Answer}}
\end{tcolorbox}
\end{flushright}
}

\newcommand{\partanswerbox}[1]{%
\begin{flushright}
\begin{tcolorbox}[
    width=0.36\textwidth,
    height=2cm,
    colback=white,
    colframe=uwpurple,
    boxrule=0.7pt,
    arc=1mm
]
{\small\textbf{Final Answer for Part #1}}
\end{tcolorbox}
\end{flushright}
}

\begin{document}

\begin{titlepage}
\centering
\vspace*{1.2cm}

{\Large \textbf{\university}}\\
{\large \department}

\vspace{1.5cm}

{\Huge \color{uwpurple}\textbf{Homework 2}}\\[0.25cm]
{\Large \textbf{\course}}\\[0.25cm]
{\large Instructor: Arman Jahangiri}

\vspace{1cm}

\begin{tcolorbox}[
    width=0.78\textwidth,
    colback=softpurple,
    colframe=uwpurple,
    boxrule=0.9pt,
    arc=3mm
]
\centering
\textbf{Submission:} Single PDF on Gradescope

\textbf{Deadline:} 10:00 PM Pacific Time on the listed due date
\end{tcolorbox}

\vspace{0.8cm}

\begin{tcolorbox}[
    width=0.86\textwidth,
    colback=white,
    colframe=uwgold,
    boxrule=0.8pt,
    arc=3mm
]
This homework is worth \textbf{50 points}.

Across the quarter, there are \textbf{eight homework assignments}, worth \textbf{400 points total}. Homework assignments together account for \textbf{40\%} of the final course grade. Further course policies can be seen in the \href{https://armanjg.github.io/MATH-394-Probability-1/}{\textbf{MATH 394 syllabus}.}
\end{tcolorbox}

\vfill
\end{titlepage}

\section*{Homework Policy}

\subsection*{Submission and Deadline}

Please:

\begin{itemize}[leftmargin=*]
    \item Submit homework as a single PDF on Gradescope by \textbf{10:00 PM (Pacific Time)} on the listed due date.
    \item Begin each problem on a new page, and ensure that pages are correctly tagged in Gradescope.
    \item Write solutions clearly and legibly.
    \item Typesetting solutions in LaTeX is encouraged, but not required.
\end{itemize}

\subsection*{Late Days}

Each student is allotted \textbf{six late days} for the quarter. A late day extends a homework deadline by up to 24 hours without penalty. For example, submitting an assignment anytime between 10:01 PM on the due date and 10:00 PM the following day counts as one late day.

The following rules apply:
\begin{itemize}[leftmargin=*]
   \item No explanation is required to use late days.
   \item At most \textbf{two late days} may be used on any single homework assignment.
   \item Late days may not be used on the final homework assignment or on any assignment for which solutions have already been released.
   \item It is the student's responsibility to keep track of how many late days have been used during the quarter.
\end{itemize}

Once all late days have been exhausted, additional late submissions will incur a penalty of \textbf{10\% per day}, up to a maximum deduction of \textbf{50\%}. Assignments submitted more than five days late, or after solutions have been released, will not be accepted.

\subsection*{Submission Issues and Technical Difficulties}

If a serious technical issue prevents a timely Gradescope submission, students may temporarily submit their assignment by email to the instructor at \href{mailto:armanjg@uw.edu}{armanjg@uw.edu}. In such cases:

\begin{itemize}[leftmargin=*]
    \item The submission must consist of a \textbf{single PDF file}. The email subject line must follow the format:
    \[
    \texttt{HW\# Submission - FULL NAME - STUDENT ID}.
    \]
    \item Once the Gradescope issue has been resolved, the student must upload the \emph{exact same file} to Gradescope as soon as possible.
    \item The student should clearly indicate in the Gradescope submission that the assignment was previously submitted by email before the deadline.
\end{itemize}

Email submissions are intended only for genuine technical emergencies and should not be used as a substitute for timely Gradescope submission.

\newpage

% ------------------------------------------------------------

\begin{problemheader}
\textbf{Problem 1. }
\end{problemheader}
\source{8}

\begin{enumerate}[label=(\alph*), leftmargin=*]
    \item How many ways are there to split a dozen people into 3 teams, where one team has 2 people, and the other two teams have 5 people each?

    \item How many ways are there to split a dozen people into 3 teams, where each team has 4 people?
\end{enumerate}

\vfill
 
\newpage

% ------------------------------------------------------------

\begin{problemheader}
\textbf{Problem 2.}
\end{problemheader}
\source{20}

Show using a story proof that
\[
{ k \choose k}
+
{ k+1 \choose k}
+
{ k+2 \choose k}
+
\cdots
+
{ n \choose k}
=
{ n+1 \choose k+1},
\]
where $n$ and $k$ are positive integers with $n\ge k$. This is called the \emph{hockey stick identity}.

\emph{Hint: Imagine arranging a group of people by age, and then think about the oldest person in a chosen subgroup.}

\vfill
 
\newpage

% ------------------------------------------------------------

\begin{problemheader}
\textbf{Problem 3. }
\end{problemheader}
\source{43}

Show that for any events $A$ and $B$,
\[
P(A)+P(B)-1
\le
P(A\cap B)
\le
P(A\cup B)
\le
P(A)+P(B).
\]

For each of these three inequalities, give a simple criterion for when the inequality is actually an equality. For example, give a simple condition such that $P(A\cap B)=P(A\cup B)$ if and only if the condition holds.

\vfill
 
 \newpage

% ------------------------------------------------------------

\begin{problemheader}
\textbf{Problem 4. }
\end{problemheader}
\source{44}

Let $A$ and $B$ be events. The \emph{difference} $B-A$ is defined to be the set of all elements of $B$ that are not in $A$. Show that if $A\subseteq B$, then
\[
P(B-A)=P(B)-P(A),
\]
directly using the axioms of probability.

\vfill
 
 \newpage

% ------------------------------------------------------------

\begin{problemheader}
\textbf{Problem 5. }
\end{problemheader}
\source{50}

A card player is dealt a 13-card hand from a well-shuffled, standard deck of cards.

What is the probability that the hand is void in at least one suit? Here, ``void in a suit'' means having no cards of that suit.

\vfill
 \newpage

% ------------------------------------------------------------

\begin{problemheader}
\textbf{Problem 6.}
\end{problemheader}
\source{55}

A club consists of 10 seniors, 12 juniors, and 15 sophomores. An organizing committee of size 5 is chosen randomly, with all subsets of size 5 equally likely.

\begin{enumerate}[label=(\alph*), leftmargin=*]
    \item Find the probability that there are exactly 3 sophomores in the committee.

    \item Find the probability that the committee has at least one representative from each of the senior, junior, and sophomore classes.
\end{enumerate}

\vfill
 

\end{document}